% Generated draft: requires mathematical and source review.
\documentclass[11pt]{article}
\usepackage[margin=25mm]{geometry}
\usepackage{fontspec}
\setmainfont{Latin Modern Roman}
\usepackage{amsmath,amssymb,mathtools,unicode-math}
\setmathfont{Latin Modern Math}
\usepackage{enumitem,needspace,xurl,hyperref}
\setlength{\parindent}{0pt}
\setlength{\parskip}{5pt}
\setlength{\emergencystretch}{3em}
\newcommand{\N}{\mathbb{N}}
\newcommand{\Z}{\mathbb{Z}}
\newcommand{\Q}{\mathbb{Q}}
\newcommand{\R}{\mathbb{R}}
\newcommand{\C}{\mathbb{C}}
\newcommand{\F}{\mathbb{F}}
\newcommand{\A}{\mathbb{A}}
\newcommand{\PP}{\mathbb{P}}
\newcommand{\E}{\mathbb{E}}
\newcommand{\eps}{\varepsilon}
\newcommand{\dd}{\,\mathrm{d}}
\begin{document}
\section*{UnsolvedMath ID: 30002000}
\textbf{Source problem code:} catalogue entry
\subsection{Definitions and mathematical statement}
Let \(G\) be a semisimple simply connected group. The supplied record does not specify its base field. Let \(W\) denote its Weyl group. Write \(\Lambda\) for the weight lattice intended by the source and \(\Lambda_+\) for its dominant cone. The notation \(\check\lambda\) is retained from the source; the record does not resolve whether these are weights of \(G\), coweights of \(G\), or weights of a dual group, nor does it specify the choice determining dominance.

Write \(\Psi(q,z)=(\Psi_{\check\lambda}(q,z))_{\check\lambda\in\Lambda}\) for a family of the intended \(q\)-Whittaker polynomials. Let \(\mathcal P\) stand for their intended coefficient and function space, and let \(\mathcal A\) stand for the intended algebra of functions in the spectral variable \(z\), with its intended \(W\)-action. These are placeholders for data absent from the supplied record. Likewise, \(M_f^{\mathrm{lat}}\) denotes the source's lattice quantum Toda operator associated with \(f\); its construction, normalization, and action on lattice-indexed families are not supplied.

\textbf{Source-level conjecture schema.} For every semisimple simply connected group \(G\), with the intended source conventions fixed, there exists exactly one family \(\Psi\in\mathcal P^{\Lambda}\) such that
\[
\begin{aligned}
\Psi_{\check\lambda}(q,z)&=0 &&\text{for every }\check\lambda\in\Lambda\setminus\Lambda_+,\\
\Psi_0(q,z)&=1,\\
M_f^{\mathrm{lat}}\bigl(\Psi(q,z)\bigr)&=f(z)\Psi(q,z)
&&\text{for every }f\in\mathcal A^W.
\end{aligned}
\]
The last equality is an equality of lattice-indexed families, with multiplication by \(f(z)\) on each component. The source does not specify the domain of \(z\), the coefficient ring, whether \(q\) is formal or specialized, or any restrictions on its values. Consequently, this schema preserves the supplied existence and uniqueness claim but is not a fully specified mathematical conjecture until those data and the operators are recovered. No replacement by a modified or related Toda formulation is intended.
\subsection{Short English statement}
For each semisimple simply connected group, the proposed family of \(q\)-Whittaker polynomials should be uniquely determined by vanishing at nondominant indices, normalization at zero, and simultaneous lattice quantum Toda eigenvalue equations for every Weyl-invariant function in the intended algebra. The supplied record leaves the lattice, function spaces, parameter domains, and operators insufficiently defined, so the formulation requires source verification and human mathematical review.
\subsection{Sources}
\begin{itemize}
\item Local UnsolvedMath registry, permanent ID \(30002000\), exact source problem code \texttt{catalogue entry}, title: \emph{\(q\)-Whittaker Eigenfunctions of Lattice Quantum Toda Operators}. The supplied record belongs to the Oberwolfach Reports Open Problems collection and gives proposed year 2012; no individual proposer is identified.
\item Supplied primary report reference: \emph{Enveloping Algebras and Geometric Representation Theory} (2012), Oberwolfach Reports, DOI 10.4171/owr/2012/13, pp. 733--809. Literature-source label: Oberwolfach Report 13 2012. \url{https://doi.org/10.4171/owr/2012/13}. The record describes these reports as documenting workshops and their open-problem discussions; it supplies no narrower locator for this conjecture.
\item Supplied literature reference: \emph{Weyl Modules and q-Whittaker Functions}. \url{https://arxiv.org/abs/1203.1583}.
\item Literature assessment dated exactly \texttt{2026-08-21}: the supplied curation labels the problem open and states that the original lattice quantum Toda joint-eigenfunction conjecture is explicitly presented as unproved, while later work proves modified or related formulations rather than this exact \(M_f^{\mathrm{lat}}\) statement in full generality. This is supplied dated curation derived from the OWR record, not an independently verified literature assessment. The record expressly requests independent verification of time-sensitive status and cited claims.
\end{itemize}
\end{document}
